\section{历史发展}\label{sec:history}

庞加莱猜想的历史是一部“问题先于技术”的历史：一个 1904 年的组合拓扑问题，等待了几何分析与偏微分方程工具的成熟，才在百年后得到解答。本节按四个阶段梳理这一历程。

\subsection{起源：庞加莱的两次出错与最终提问（1895--1904）}

庞加莱在《位置分析》（\textit{Analysis Situs}）及其五篇补充篇\cite{poincare1895,poincare1904}中创建了代数拓扑的基本工具。在第二补充篇（1900）中，他曾断言：若两个流形有相同的贝蒂数与挠率系数（即同调群同构），则它们（在某些条件下）同胚。不久他本人发现了反例——后来称为\textbf{庞加莱同调球面}的十二面体空间：它由正十二面体的对面粘合而成，与 $S^3$ 有完全相同的同调，但基本群是 120 阶的非平凡群。这一发现迫使他引入基本群，并意识到同调的不足。

在第五补充篇\cite{poincare1904}中，庞加莱提出：“是否可以认为，基本群平凡的紧流形恰好是（三维）球面？”——他没有称之为定理或猜想，只是“尚不能确定的一个问题”。这个谦逊的问号成为此后一百年三维拓扑的中心问题。值得注意的是，庞加莱同时代的工具（同调、基本群、覆叠空间）已经穷尽：任何更进一步的组合不变量在半个多世纪后才出现。

\subsection{早期攻坚与三维技术的积累（1930--1980）}

二十世纪上半叶的尝试屡屡失败，但失败孕育了三维拓扑本身：

\begin{itemize}
    \item \textbf{怀特海德（1934）}宣布了证明\cite{whitehead1934}，随后发现漏洞；检查漏洞的过程中他构造了著名的\textbf{怀特海流形}——单连通、可缩却不与 $\mathbb{R}^3$ 同胚的开三维流形，说明三维中“单连通”与“无洞”之间隔着真正的鸿沟。
    \item \textbf{Dehn 引理与圈定理}：Dehn（1912）的引理存在漏洞，Papakyriakopoulou 在 1957 年给出第一个完整证明\cite{papakyriakou1957}，随后又证明了圈定理与球面定理。这组“三维拓扑三大定理”把许多问题约化到较简单的流形，成为此后一切工作的地基。
    \item \textbf{Haken--Waldhausen 理论}：Haken 引入不可压缩曲面与层次分解\cite{haken1962}，Waldhausen 证明大类 Haken 流形的刚性（同调或群论信息决定拓扑）。Haken 流形不可能是庞加莱猜想的反例，这使得候选反例被逐步排除。
    \item \textbf{其他贡献}：Bing 的收缩理论证明了“双重cover 定理”等判据；Fox 与 Artin 编织了基本群计算的技术；Milnor（1965）对群中 Dehn 函数的分析暗示了“胖”群难以成为三维流形群。这些工作逐步把猜想的可能反例压缩到越来越窄的角落。
\end{itemize}

\subsection{高维的胜利与三维的孤独（1960--1982）}

出人意料的是，难度曲线在维数上并不单调：

\begin{itemize}
    \item \textbf{$n \geq 5$}：斯梅尔借助 $h$-协边界定理\cite{smale1961}证明了高维广义庞加莱猜想（斯梅尔因此获 1966 年菲尔兹奖）；Stallings 与 Zeeman 独立给出了部分证明。其技术核心——把流形分解为把柄并整理附件次序——本质上依赖于高维中把柄有足够的空间摆布。
    \item \textbf{$n = 4$（拓扑）}：四维中把柄附件失去自由度，h-协边界定理失效。弗里德曼以全新的圆盘嵌入定理攻克了拓扑四维情形\cite{freedman1982}，并据此分类了单连通闭四维流形的交型。几乎同时，Donaldson 用规范场论揭示了光滑四维结构的病态\cite{donaldson1983}：光滑范畴中存在拓扑方法不可见的约束。两者结合立刻推出 $\mathbb{R}^4$ 存在怪异光滑结构——而这正暗示\textbf{光滑四维庞加莱猜想处于完全不同的难度层级}。
    \item \textbf{$n = 3$}：把柄理论在此退化（把柄与骨架纠缠），组合方法止步。三维必须等待新的语言。
\end{itemize}

\subsection{几何化的转向与里奇流（1970--2002）}

真正的转折来自对“三维流形长什么样”的重新想象：

\begin{itemize}
    \item \textbf{Thurston 几何化猜想（1982）}\cite{thurston1982}：任何闭三维流形都可沿球面与环面特征子流形分解为若干块，每块 admit 八种模型几何之一（$S^3$、$\mathbb{E}^3$、$H^3$、$S^2 \times \mathbb{R}$、$H^2 \times \mathbb{R}$、$\widetilde{\mathrm{SL}_2\mathbb{R}}$、Nil、Sol）。单连通性排除了除 $S^3$ 外的一切几何，庞加莱猜想成为几何化的直接推论。Thurston 本人为 Haken 流形证明了几何化（并获得 1982 年菲尔兹奖），但双曲化定理之外的推广（尤其是含 $S^3$ 块的情形）停滞不前。
    \item \textbf{Hamilton 引入里奇流（1982）}\cite{hamilton1982}：$\partial_t g = -2\operatorname{Ric}$ 是曲率的扩散方程，如同热方程把温度均匀化，它把任意度量“拉伸”向常曲率度量。Hamilton 证明了二维情形的收敛与三维的正 Ricci 曲率情形（直接得到 $S^3$ 球定理）。但在一般三维流形上，流会在有限时间产生奇点——瓶颈与尖点——程序中断。哈密顿纲领（1990 年代）：理解所有奇点，切除并继续，最终得到几何分解——成为此后十年的路线图，其关键部分（Harnack 估计、球定理、某些奇点分析）由 Hamilton 与 Ivey 等人完成，但一般奇点的分类卡住了。
\end{itemize}

历史在 2002 年 11 月转折：格里戈里·佩雷尔曼在 arXiv 贴出第一篇预印本。三篇论文\cite{perelman2002entropy,perelman2003surgery,perelman2003extinction}在两年内补齐了哈密顿纲领缺失的所有环节，这个故事的细节见\cref{sec:results}。
